\documentclass[11pt]{article}

%------------------------------------------------
% PAQUETES
%------------------------------------------------
\usepackage[spanish]{babel}
\usepackage[utf8]{inputenc}
\usepackage[T1]{fontenc}
\usepackage{lmodern}

\usepackage{amsmath,amssymb}
\usepackage[a4paper,margin=2.5cm]{geometry}
\usepackage{xcolor}
\usepackage{graphicx}
\usepackage{fancyhdr}
\usepackage[most]{tcolorbox}
\usepackage{enumitem}

%------------------------------------------------
% COLORES
%------------------------------------------------
\definecolor{Titulo}{HTML}{0F2D52}
\definecolor{Subtitulo}{HTML}{4E5D6C}
\definecolor{Caja}{HTML}{F5F7FA}
\definecolor{Borde}{HTML}{D6DEE6}

%------------------------------------------------
% ENCABEZADO Y PIE
%------------------------------------------------
\pagestyle{fancy}
\fancyhf{}
\lhead{\textcolor{Subtitulo}{Ayudantía 1}}
\rhead{\textcolor{Subtitulo}{Lógica proposicional}}
\cfoot{\thepage}
\renewcommand{\headrulewidth}{0.4pt}

%------------------------------------------------
% CAJAS
%------------------------------------------------
\newtcolorbox{ejercicio}{
	colback=white,
	colframe=Borde,
	boxrule=0.7pt,
	arc=3pt,
	left=6pt,
	right=6pt,
	top=6pt,
	bottom=6pt
}

\newtcolorbox{solucion}{
	colback=Caja,
	colframe=Borde,
	boxrule=0.7pt,
	arc=3pt,
	left=6pt,
	right=6pt,
	top=6pt,
	bottom=6pt
}

%------------------------------------------------
% DOCUMENTO
%------------------------------------------------
\begin{document}
	
	\begin{center}
		% \includegraphics[width=2cm]{logo_ust.png}
		
		{\LARGE\bfseries\textcolor{Titulo}{Ayudantía 1}}
		
		\vspace{3mm}
		
		{\large\bfseries\textcolor{Subtitulo}{Lógica proposicional}}
		
		\vspace{4mm}
		
		{\normalsize Manuel Concha}
		
		{\small Universidad Santo Tomás}
	\end{center}
	
	\vspace{8mm}
	
	{\Large\bfseries\textcolor{Titulo}{Ejercicios}}
	
	\vspace{4mm}
	
	%------------------------------------------------
	% EJERCICIO 1
	%------------------------------------------------
	\begin{ejercicio}
		\textbf{Ejercicio 1.} Indique cuáles de los siguientes enunciados son proposiciones.
		
		\begin{enumerate}[label=\arabic*.]
			\item El precio del producto aumenta.
			\item ¿Cuál es el valor del dólar hoy?
			\item 15 es divisible por 3.
			\item ¡Excelente inversión!
		\end{enumerate}
	\end{ejercicio}
	
	\begin{solucion}
		\textbf{Solución}
		
		\begin{itemize}
			\item (1) Sí es proposición.
			\item (2) No es proposición, porque es una pregunta.
			\item (3) Sí es proposición y es verdadera.
			\item (4) No es proposición, porque es una exclamación.
		\end{itemize}
	\end{solucion}
	
	\vspace{5mm}
	
	%------------------------------------------------
	% EJERCICIO 2
	%------------------------------------------------
	\begin{ejercicio}
		\textbf{Ejercicio 2.} Sean las proposiciones:
		\[
		p:\ \text{``Aumenta el precio''}
		\qquad
		q:\ \text{``Disminuye la demanda''}
		\]
		
		Escriba en lenguaje usual:
		
		\begin{enumerate}[label=\alph*)]
			\item \(p \land q\)
			\item \(p \lor q\)
			\item \(\neg p\)
		\end{enumerate}
	\end{ejercicio}
	
	\begin{solucion}
		\textbf{Solución}
		
		\begin{itemize}
			\item \(p \land q\): Aumenta el precio y disminuye la demanda.
			\item \(p \lor q\): Aumenta el precio o disminuye la demanda.
			\item \(\neg p\): No aumenta el precio.
		\end{itemize}
	\end{solucion}
	
	\vspace{5mm}
	
	%------------------------------------------------
	% EJERCICIO 3
	%------------------------------------------------
	\begin{ejercicio}
		\textbf{Ejercicio 3.} Considere las proposiciones:
		\[
		p:\ \text{``8 es par''}
		\qquad
		q:\ \text{``8 es múltiplo de 3''}
		\]
		
		Determine el valor de verdad de:
		
		\begin{enumerate}[label=\arabic*.]
			\item \(p\)
			\item \(q\)
			\item \(p \land q\)
			\item \(p \lor q\)
		\end{enumerate}
	\end{ejercicio}
	
	\begin{solucion}
		\textbf{Solución}
		
		\[
		p = V
		\qquad
		q = F
		\]
		
		\begin{itemize}
			\item \(p\) es verdadero.
			\item \(q\) es falso.
			\item \(p \land q = F\).
			\item \(p \lor q = V\).
		\end{itemize}
	\end{solucion}
	
\end{document}